\chapter{Trabajo a futuro y Conclusiones}

Durante TT2 se continuará con la fase de desarrollo e implementación del proyecto, tomando como base la investigación, documentación y análisis realizados previamente.

Como parte del trabajo futuro, se desarrollará la integración del sistema, basado en diagramas de arquitectura, diseño de módulos, flujo de procesamiento de datos e integración para la interacción entre componentes de hardware y software.

Se realizará la implementación progresiva de los módulos previamente investigados, integrando sistemas de visión computacional, reconocimiento de voz, navegación y generación de instrucciones auditivas dentro de una plataforma funcional.

Asimismo, se llevará a cabo el entrenamiento y ajuste de los modelos de inteligencia artificial mediante pruebas controladas y recopilación de datos en escenarios representativos, permitiendo evaluar el comportamiento del sistema ante distintos tipos de obstáculos y condiciones ambientales.

También se documentarán las pruebas experimentales realizadas al prototipo, registrando métricas de desempeño relacionadas con precisión de detección, tiempos de respuesta, estabilidad operativa y efectividad de navegación.

Durante esta etapa se buscará validar el funcionamiento integral del sistema mediante pruebas en entornos interiores y exteriores, analizando la capacidad del prototipo para asistir al usuario durante recorridos básicos de navegación.

Adicionalmente, se contempla continuar investigando estrategias de optimización orientadas a mejorar la eficiencia del procesamiento en hardware embebido, así como técnicas de navegación adaptativa y recalculación de rutas en tiempo real.

Como objetivo principal, se busca obtener un prototipo funcional y documentado que permita validar la propuesta tecnológica planteada en la etapa de investigación, demostrando la factibilidad del sistema como herramienta de asistencia para personas con visión reducida.

\section{Conclusiones de Trabajo Terminal I}
A partir de la investigación y análisis realizados durante TT1, se determinó la viabilidad técnica del desarrollo de un sistema de navegación asistida para personas con visión reducida, integrando tecnologías de inteligencia artificial, visión computacional y procesamiento de voz.

La etapa de investigación permitió identificar los principales desafíos asociados a la movilidad autónoma asistida, así como las tecnologías más adecuadas para abordar problemas relacionados con detección de obstáculos, orientación espacial y emisión de instrucciones auditivas en tiempo real.

Se estableció una arquitectura modular para el sistema, permitiendo definir de manera estructurada cada uno de los componentes que conformarán el prototipo funcional en TT2. Esta arquitectura facilita la documentación, validación e integración progresiva de los módulos de hardware y software.

Asimismo, se definieron los requerimientos funcionales y no funcionales del sistema, considerando aspectos fundamentales como precisión, latencia, autonomía operativa y capacidad de procesamiento en una raspberry pi.

La investigación tecnológica realizada permitió justificar el uso de herramientas como YOLOv8 para detección de objetos, Vosk para reconocimiento de voz, spaCy para procesamiento de lenguaje natural y algoritmos de navegación basados en SLAM y aprendizaje profundo por refuerzo, debido a su compatibilidad con entornos de tiempo real y hardware de recursos limitados.

De igual manera, se documentaron los flujos operativos, casos de uso, interacciones del sistema y posibles escenarios de funcionamiento, estableciendo una base sólida para el desarrollo e implementación del prototipo durante la siguiente etapa del proyecto terminal.

Finalmente, el proyecto se sustenta en una problemática social relevante, orientada a generar herramientas tecnológicas de apoyo que contribuyan a mejorar la autonomía y seguridad de las personas con discapacidad visual durante sus desplazamientos.